Agentes baseados em grandes modelos de linguagem podem interpretar documentos e acionar ferramentas corporativas, mas também podem tratar instruções em conteúdo externo como comandos legítimos. Em Contas a Pagar, uma injeção de prompt indireta pode alterar campos financeiros, redirecionar o beneficiário ou induzir ação crítica sem autorização. Este trabalho objetiva desenvolver e avaliar uma arquitetura de proteção para integrar LLMs ao Odoo 19.0 Community em um fluxo sintético de contas de fornecedor. A pesquisa segue Design Science Research com experimento controlado, comparando uma integração funcional de referência (C0) à defesa em profundidade (C1), com política determinística, menor privilégio, aprovação vinculada, idempotência e auditoria. O experimento principal prevê 90 pares contrafactuais, três repetições e 1.080 execuções com um modelo DeepSeek; o piloto separado acrescenta 144. Há provas técnicas do fluxo documental, dos caminhos C0/C1, da captura e restauração do estado e de registros que preservam falhas. No piloto de propostas, em duas rodadas, a segunda com estratégias de ataque revisadas, nenhuma das 72 propostas adversariais de C0 incorporou o objetivo do atacante, e as propostas benignas foram exatas em 36 de 36 execuções em C0 e 35 de 36 em C1. Com esse modelo e esses ataques estáticos, a redução do sucesso dos ataques não pode ser demonstrada; um enquadramento estimativo aguarda decisão do grupo. O efeito no estado do ERP e o teste final permanecem pendentes. Replicação entre modelos, ablação e ataques adaptados à defesa ficam como trabalhos futuros.
